\chapter{Limites, validation et responsabilité}

\section{Limites actuelles}

Les limites exactes dépendent de la version du logiciel et du moteur choisi. À
ce stade, il faut notamment vérifier :

\begin{itemize}
  \item les formulations de plaques disponibles ;
  \item les conventions de signe des diagrammes ;
  \item les résultats multi-cas et enveloppes ;
  \item les vérifications normatives encore en développement ;
  \item les cas avancés non linéaires, sismiques ou temporels.
\end{itemize}

\section{Validation recommandée}

Pour chaque nouveau type de modèle, il est recommandé de comparer les résultats
avec :

\begin{itemize}
  \item une solution analytique simple ;
  \item un autre logiciel de calcul reconnu ;
  \item des cas tests internes versionnés ;
  \item les ordres de grandeur attendus par l'expérience métier.
\end{itemize}

\section{Avertissement}

\begin{warningbox}
\logiciel{} est un outil d'aide à la modélisation et au calcul. L'utilisateur
reste seul responsable de l'interprétation des résultats, de la validation du
modèle, du choix des hypothèses et de l'utilisation professionnelle des sorties.
\end{warningbox}

\section{Traçabilité}

Pour une utilisation sérieuse, conserver :

\begin{itemize}
  \item la version du logiciel ;
  \item le fichier projet ;
  \item les hypothèses de charges ;
  \item les conventions de signe retenues ;
  \item les captures ou exports de résultats ;
  \item les vérifications indépendantes.
\end{itemize}
